% ============================================================
%  Formelsammlung  –  Systemtheorie 1
%  Querformat, 3-spaltig, thematisch gruppiert
% ============================================================
\clearpage
\newgeometry{landscape, top=1.1cm, bottom=1.1cm, left=1.1cm, right=1.1cm}
\pagestyle{empty}
\addcontentsline{toc}{chapter}{Formelsammlung}

% ── kompaktes Layout ─────────────────────────────────────────
\setlength{\parindent}{0pt}
\setlength{\abovedisplayskip}{2pt}
\setlength{\belowdisplayskip}{2pt}
\setlength{\abovedisplayshortskip}{1pt}
\setlength{\belowdisplayshortskip}{1pt}
\setlength{\columnsep}{1.4em}
\setlength{\columnseprule}{0.4pt}

% Themen-Kopf: Titel + kleiner Skript-Verweis + Trennlinie
\newcommand{\ft}[2]{%
  \par\smallskip\penalty-200
  {\bfseries\footnotesize #1}\nobreak\hfill{\tiny\hyperref[#2]{\S\,\ref*{#2}}}\par\nobreak
  \vspace{-0.35em}{\color{accent}\hrule height 0.4pt}\nobreak\vspace{0.2em}\nobreak}

\begin{center}
  {\color{accent}\rule{\linewidth}{1pt}}\\[0.2em]
  {\Large\bfseries Formelsammlung}\quad\small Systemtheorie\,1 -- Signale und Systeme\,I\\[0.1em]
  {\color{accent}\rule{\linewidth}{1pt}}
\end{center}
\vspace{0.4em}

\footnotesize
\begin{multicols}{3}

% ─────────────────────────────────────────────────────────────
\ft{Elementare analoge Signale}{sec:elementare-analoge-signale}
Gleichsignal $x(t)=A$.\quad Sinus:
\[ x(t)=A\sin(\omega t+\varphi)=A\cos\!\big(\omega t+\varphi-\tfrac{\pi}{2}\big) \]
\[ \omega=2\pi f=\tfrac{2\pi}{T},\quad f=\tfrac{1}{T} \]
Sprung $\varepsilon(t)=\begin{cases}0&t<0\\1&t\ge0\end{cases}$\;
Rechteck $\operatorname{rect}\!\big(\tfrac{t}{T}\big)$, Fläche $=T$.\\
Exponential $x(t)=A\,e^{-t/\tau}\varepsilon(t)$, Steigung $|_{0^+}=-\tfrac{A}{\tau}$.\\
$\operatorname{si}\!\big(\tfrac{t}{T}\big)=\dfrac{\sin(\pi t/T)}{\pi t/T}$.\quad
Dirac (Siebeigenschaft):
\[ \int_{-\infty}^{\infty}\! f(\xi)\,\delta(\xi-a)\,\mathrm d\xi=f(a),\ \int\delta\,\mathrm d\xi=1 \]

% ─────────────────────────────────────────────────────────────
\ft{Elementare diskrete Signale}{sec:elementare-diskrete-signale}
$\delta[k]=\begin{cases}1&k=0\\0&\text{sonst}\end{cases}$\quad
$\varepsilon[k]=\begin{cases}1&k\ge0\\0&k<0\end{cases}$\\
Rechteck $x[k]=\varepsilon[k]-\varepsilon[k-M]$, Energie $E=MA^2$.\\
Exp.-Folge $x[k]=a^k\varepsilon[k]$.\quad Sinus $x[k]=A\sin\!\big(2\pi k\tfrac{f}{f_A}\big)$.\\
Geometrische Reihe:
\[ \sum_{k=0}^{M-1}\! q^k=\frac{1-q^M}{1-q},\quad \sum_{k=0}^{\infty} q^k=\frac{1}{1-q}\ (|q|<1) \]

% ─────────────────────────────────────────────────────────────
\ft{Komplexe Signale \& Zeiger}{sec:komplexwertige-signale}
\[ \underline{x}(t)=x_R+jx_I=|\underline{x}|\,e^{j\phi},\quad j=\sqrt{-1} \]
\[ x_R=\operatorname{Re}\{\underline x\}=\tfrac12(\underline x+\underline x^{*}),\
   x_I=\operatorname{Im}\{\underline x\}=\tfrac{1}{2j}(\underline x-\underline x^{*}) \]
\[ |\underline x|=\sqrt{x_R^2+x_I^2},\quad \phi=\arg\{\underline x\} \]
Komplexes Exp.: $\underline x(t)=A\,e^{j(\omega t+\varphi)}=\underline A\,e^{j\omega t}$, $\underline A=A e^{j\varphi}$.

% ─────────────────────────────────────────────────────────────
\ft{Symmetrie}{sec:symmetrie-signale}
gerade $x(t)=x(-t)$;\quad ungerade $x(t)=-x(-t)$.\\
konj.\ gerade $\underline x(t)=\underline x^{*}(-t)$;\ konj.\ ungerade $\underline x(t)=-\underline x^{*}(-t)$.\\
Zerlegung $x=x_{\mathrm g}+x_{\mathrm u}$ mit
\[ x_{\mathrm g}=\tfrac12\big(x(t)+x(-t)\big),\ x_{\mathrm u}=\tfrac12\big(x(t)-x(-t)\big) \]

% ─────────────────────────────────────────────────────────────
\ft{Mittelwert, Energie, Leistung, Norm}{sec:norm-energie-leistung}
$\bar x=\dfrac{1}{t_2-t_1}\!\int_{t_1}^{t_2}\! x\,\mathrm dt=\dfrac{1}{k_2-k_1+1}\!\displaystyle\sum_{k=k_1}^{k_2}\! x[k]$;\\
periodisch $\bar x=\tfrac1T\!\int_{-T/2}^{T/2}\! x\,\mathrm dt=\tfrac1N\!\sum_{k=0}^{N-1}\! x[k]$.\\
$L_p$-Norm $\|\underline x\|_p=\big(\int|\underline x|^p\mathrm dt\big)^{1/p}$.\\
Energie $E=\int|\underline x(t)|^2\mathrm dt=\sum|\underline x[k]|^2$ ($E<\infty$: Energiesignal).\\
Leistung $P=\lim_{T\to\infty}\tfrac1T\!\int_{-T/2}^{T/2}\!|\underline x|^2\mathrm dt$;\;
periodisch $P=\tfrac1T\!\int_T\!|\underline x|^2\mathrm dt$\\
($0<P<\infty$: Leistungssignal).

% ─────────────────────────────────────────────────────────────
\ft{Inneres Produkt \& Korrelation}{chap:inneres-produkt}
$\langle \underline x,\underline y\rangle=\int \underline x^{*}(t)\underline y(t)\,\mathrm dt=\sum \underline x^{*}[k]\underline y[k]$; \ orthogonal $\Leftrightarrow \langle\cdot\rangle=0$.\\
Kreuzkorrelation (KKF):
\[ \varphi_{xy}(\tau)=\!\int_{-\infty}^{\infty}\!\underline x^{*}(t)\,\underline y(t+\tau)\,\mathrm dt \]
\[ \varphi_{xy}[\varkappa]=\sum_{k}\underline x^{*}[k]\,\underline y[k+\varkappa] \]
Autokorrelation (AKF): $\varphi_{xx}$ (KKF mit $y=x$),
$\varphi_{xx}(0)=E$, $\varphi_{xx}(\tau)=\varphi_{xx}^{*}(-\tau)$.\\
Sinus: $\tilde\varphi_{xx}(\tau)=\tfrac{A^2}{2}\cos(\omega\tau)$.

% ─────────────────────────────────────────────────────────────
\ft{Empirische Kenngrößen}{sec:streumasse}
$\hat{\bar x}=\dfrac1N\displaystyle\sum_{k=0}^{N-1}\! x[k]$;\quad
$\hat\sigma^2=\dfrac{1}{N-1}\displaystyle\sum_{k=0}^{N-1}\!\big(x[k]-\hat{\bar x}\big)^2$;\\
$\hat\sigma=\sqrt{\hat\sigma^2}$;\quad
empir.\ KKF $\hat\varphi_{xy}[\varkappa]=\sum_{k=0}^{N} x[k]\,y[k+\varkappa]$.

% ─────────────────────────────────────────────────────────────
\ft{Fourier-Reihe \& DFT}{sec:fourierreihe}
reell: $x(t)=\bar x+\sum_{n\ge1}X_a[n]\cos(n\omega_0t)+X_b[n]\sin(n\omega_0t)$.\\
komplex: $x(t)=\sum_{n=-\infty}^{\infty}\! X[n]e^{jn\omega_0t}$, $\omega_0=\tfrac{2\pi}{T_0}$,
\[ X[n]=\frac{1}{T_0}\!\int_{t_1}^{t_1+T_0}\! x(t)\,e^{-jn\omega_0t}\,\mathrm dt \]
$X_a=\tfrac{2}{T_0}\!\int x\cos$, $X_b=\tfrac{2}{T_0}\!\int x\sin$.\\
DFT:
\[ x[k]=\sum_{m=0}^{M-1}\! X[m]\,e^{j\frac{2\pi}{M}mk} \]
\[ X[m]=\tfrac1M\sum_{k=0}^{M-1}\! x[k]\,e^{-j\frac{2\pi}{M}mk} \]
Bandbreite $B=f_{\max}-f_{\min}$; \ $f_m=\tfrac{m}{M}\tfrac{1}{T_A}$.

% ─────────────────────────────────────────────────────────────
\ft{Abtastung \& Quantisierung}{sec:abtasttheorem}
$f_A=\tfrac{\text{Abtastwerte}}{\text{Zeit}}$, $T_A=\tfrac{1}{f_A}$.\\
Abtasttheorem $f_A>2f_g$; \ Aliasing/Spiegelfreq.\ $\tilde f=f_A-f$.\\
Ideale Rekonstruktion:
\[ \tilde x(t)=\sum_k x[k]\,\operatorname{si}\!\big(\tfrac{t-kT_A}{T_A}\big) \]
Halteglied $\tilde x(t)=x(kT_A)$, $kT_A\le t<(k{+}1)T_A$.\\
Quant.-Fehler $e=x-x_q$, $|e|\le\tfrac{\Delta x}{2}$, Leistung $P_e=\dfrac{\Delta x^2}{12}$.\\
Dezibel: $10\lg\!\tfrac{P_2}{P_1}$ [dB] $=20\lg\!\tfrac{A_2}{A_1}$ [dB].

% ─────────────────────────────────────────────────────────────
\ft{Systeme \& Zustandsraum}{sec:zustandsraum}
System $\mathbb S=\{(x(t),y(t))\mid t_0\le t<t_1\}$.\\
kontinuierlich:
\[ \tfrac{\mathrm d\vec z}{\mathrm dt}=\mathbf A\vec z+\mathbf B\vec x,\quad \vec y=\mathbf C\vec z+\mathbf D\vec x \]
diskret:
\[ \vec z[k{+}1]=\mathbf A\vec z[k]+\mathbf B\vec x[k] \]
\[ \vec y[k]=\mathbf C\vec z[k]+\mathbf D\vec x[k] \]
\emph{LTI}: zeitinvariant $y(t-\tau_0)=\mathcal S[x(t-\tau_0)]$ und linear
$a y_1+b y_2=\mathcal S[a x_1+b x_2]$.\\
kausal: $y(t_0)$ nur von $t\le t_0$;\ BIBO: $|x|<\infty\Rightarrow|y|<\infty$.

% ─────────────────────────────────────────────────────────────
\ft{LTI-Systeme \& Faltung}{sec:faltungsintegral}
Impulsantwort $h=\mathcal S[\delta]$. Faltung:
\[ y(t)=\!\int_{-\infty}^{\infty}\! x(\xi)\,h(t-\xi)\,\mathrm d\xi=x*h \]
\[ y[k]=\sum_{l} x[l]\,h[k-l]=x[k]*h[k] \]
$x*\delta=x$;\ $x*h=h*x$;\ assoziativ;\ distributiv.\\
kausal $h(t)=0\ \forall t<0$;\ BIBO $\int|h|\mathrm dt<\infty$ ($\sum|h[l]|<\infty$).\\
Frequenzgang $\underline H(j\omega)=\int h(\xi)e^{-j\omega\xi}\mathrm d\xi$ (Eigenfkt.\ $e^{j\omega t}$).\\
Wiener--Lee: $\varphi_{yy}=\varphi_{hh}*\varphi_{xx}$.

% ─────────────────────────────────────────────────────────────
\ft{Parametrische Systeme (FIR/IIR)}{sec:parametrische-lsi}
allg.\ LSI: $y[k]=\sum_{\nu=0}^{N} x[k{-}\nu]b_\nu+\sum_{\mu=1}^{M} y[k{-}\mu]a_\mu$.\\
FIR: $y[k]=\sum_{\nu=0}^{N} x[k{-}\nu]b_\nu$.\\
IIR: $y[k]=x[k]b_0+\sum_{\mu=1}^{M} y[k{-}\mu]a_\mu$.

% ─────────────────────────────────────────────────────────────
\ft{Wahrscheinlichkeit -- Grundlagen}{sec:axiome-wahrscheinlichkeit}
Axiome: $P(A)\ge0$;\ $P(\Omega)=1$;\ $P(A{+}B)=P(A)+P(B)$ (disjunkt).\\
Laplace: $P(\xi)=\tfrac{1}{|\Omega|}$.\quad
unabh.\ Versuche: $P((\xi_i,\xi_j))=P(\xi_i)P(\xi_j)$.\\
Rand aus Verbund: $P(A_i)=\sum_j P(A_i,B_j)$;\ unabh.\ $P(A,B)=P(A)P(B)$.

% ─────────────────────────────────────────────────────────────
\ft{Kombinatorik}{sec:kombinatorik}
\renewcommand{\arraystretch}{1.9}
\begin{tabular}{@{}p{1.5cm}cc@{}}
  \toprule
   & \scriptsize o.\ Zurückl. & \scriptsize m.\ Zurückl. \\
  \midrule
  \scriptsize geordnet   & $\dfrac{n!}{(n-k)!}$ & $n^k$ \\
  \scriptsize ungeordnet & $\dbinom{n}{k}$      & $\dbinom{n{-}1{+}k}{k}$ \\
  \bottomrule
\end{tabular}
\renewcommand{\arraystretch}{1}\\[0.2em]
Permutationen (alle $n$): $n!$.

% ─────────────────────────────────────────────────────────────
\ft{Bedingte W.\ \& Satz von Bayes}{sec:satz-von-bayes}
\[ P(A_i|B_j)=\frac{P(A_i,B_j)}{P(B_j)} \]
totale W.: $P(A_k)=\sum_{i} P(A_k|B_i)\,P(B_i)$.
\[ P(B_j|A_k)=\frac{P(A_k|B_j)P(B_j)}{\sum_{i} P(A_k|B_i)P(B_i)} \]

% ─────────────────────────────────────────────────────────────
\ft{Zufallsvariablen \& Momente}{sec:momente}
diskret: $\mathrm E\{X\}=\sum_i x_i p_X(x_i)$, $\mathrm E\{X^2\}=\sum_i x_i^2 p_X(x_i)$.\\
kontinuierlich: $\mathrm E\{g(X)\}=\int g(x)p_X(x)\,\mathrm dx$.
\[ \operatorname{var}\{X\}=\mathrm E\{(X-\mu)^2\}=\mathrm E\{X^2\}-\mathrm E\{X\}^2=\sigma^2 \]
$\sigma=\sqrt{\operatorname{var}\{X\}}$;\\
$\mathrm E\{aX{+}bY\}=a\,\mathrm E\{X\}+b\,\mathrm E\{Y\}$.

% ─────────────────────────────────────────────────────────────
\ft{Diskrete Verteilungen}{sec:binomialverteilung}
\emph{Binomial:} $P(k)=\dbinom{N}{k}P^k(1-P)^{N-k}$,
\[ \mathrm E\{X\}=NP,\qquad \operatorname{var}\{X\}=NP(1-P) \]
\emph{Poisson:} $p(k)=e^{-\lambda}\dfrac{\lambda^k}{k!}$,\;
$\mathrm E\{X\}=\lambda$, $\mathrm E\{X^2\}=\lambda^2+\lambda$, $\operatorname{var}=\lambda$.

% ─────────────────────────────────────────────────────────────
\ft{Stetige Verteilungen}{sec:normalverteilung}
{\scriptsize
\renewcommand{\arraystretch}{2.1}
\begin{tabular}{@{}p{1.15cm}ccc@{}}
  \toprule
  Vert. & $p_X(x)$ & $\mathrm E\{X\}$ & $\operatorname{var}\{X\}$ \\
  \midrule
  Gleich $[a,b]$ & $\tfrac{1}{b-a}$ & $\tfrac{a+b}{2}$ & $\tfrac{(b-a)^2}{12}$ \\
  Normal & $\tfrac{1}{\sqrt{2\pi\sigma^2}}e^{-\frac{(x-\mu)^2}{2\sigma^2}}$ & $\mu$ & $\sigma^2$ \\
  Exp.\ ($x{\ge}0$) & $\lambda e^{-\lambda x}$ & $\tfrac{1}{\lambda}$ & $\tfrac{1}{\lambda^2}$ \\
  Laplace & $\tfrac{1}{2\sigma}e^{-\frac{|x-\mu|}{\sigma}}$ & $\mu$ & $2\sigma^2$ \\
  \bottomrule
\end{tabular}
\renewcommand{\arraystretch}{1}}

% ─────────────────────────────────────────────────────────────
\ft{Verteilungsfkt.\ \& Transformation}{sec:transformation-zv}
$F_X(x)=P(X\le x)=\int_{-\infty}^{x} p_X(u)\,\mathrm du$,\ $p_X=\tfrac{\mathrm d}{\mathrm dx}F_X$.\\
monoton $Y=g(X)$: $p_Y(y)=\dfrac{p_X(g^{-1}(y))}{|g'(g^{-1}(y))|}$.\\
nichtmonoton: $p_Y(y)=\sum_i \dfrac{p_X(x_i)}{|g'(x_i)|}$.

% ─────────────────────────────────────────────────────────────
\ft{Zwei Zufallsvariablen}{sec:kovarianz-korrelation}
$F_{XY}=P(X\le x,Y\le y)$,\ $p_{XY}=\dfrac{\partial^2 F_{XY}}{\partial x\,\partial y}$.\\
Rand $p_X=\int p_{XY}\,\mathrm dy$;\ unabhängig $p_{XY}=p_X p_Y$.\\
Kovarianz $C=\mathrm E\{(X-\mu_X)(Y-\mu_Y)\}$;\\
$r_{XY}=\dfrac{C}{\sigma_X\sigma_Y}$, $|r_{XY}|\le1$;\ orthogonal $\mathrm E\{XY\}=0$.\\
Summe unabh.\ ZV: $p_Z=p_X*p_Y$.

% ─────────────────────────────────────────────────────────────
\ft{Informationstheorie}{sec:entropie-redundanz}
Entscheidungsgehalt $H_0=\log_2 N$.\quad
Info-Gehalt $I(a_i)=-\log_2 P(a_i)$.
\[ H=-\sum_{i} P(a_i)\log_2 P(a_i)\ \le H_0 \]
rel.\ Entropie $\tfrac{H}{H_0}$;\ Redundanz $r_q=1-\tfrac{H}{H_0}$.\\
Kanalkapazität $C=\lim_{\tau\to\infty}\tfrac{\log_2 N(\tau)}{\tau}$.\\
Coderedundanz $r_c=\tfrac{\mathrm E\{L\}-H}{H}$;\
KL $D=\sum_i P\log_2\tfrac{P}{\widehat P}$.

\end{multicols}

\restoregeometry
